\documentclass[11pt,a4paper,oneside,openany]{book}
\usepackage{biblioflow-ista}

\renewcommand{\titolodocumento}{Master Test Plan}
\renewcommand{\sottotitolodocumento}{Master Test Plan}

\begin{document}
\frontmatter
\copertina
\tableofcontents

\mainmatter
\chapter{Introduzione}

Questo documento definisce le strategie pianificate per il testing di
BiblioFlow. Il piano definisce i diversi livelli di test: test di unità, di
integrazione, di sistema ed E2E, con l'obiettivo di garantire la qualità del
software considerando anche le nuove funzionalità descritte nel documento
Change Requests.

Verranno escluse dal perimetro di testing le integrazioni con servizi esterni
che il progetto usa ma non realizza: il recapito effettivo delle email, il
fornitore di accesso Google e gli strumenti di monitoraggio.

\chapter{Strategie di Testing}

\section{Testing di unità}

La fase di testing di unità è finalizzata alla verifica dei singoli componenti
di BiblioFlow, con l'intento di isolare e validare il comportamento dei moduli
critici: servizio di dominio delle prenotazioni, autenticazione, automazioni e
route handler. La metodologia adottata è di tipo white-box, focalizzata sui
criteri di statement coverage e branch coverage misurati con Vitest. Le
dipendenze esterne, cioè il client Prisma e la sessione di Auth.js, sono
sostituite da mock, così che ogni test verifichi solo la logica del modulo.

Allo stesso livello appartengono i test dei componenti React coinvolti dalle
Change Request (mappa dei posti, griglia mobile, lista d'attesa, azioni dello
staff), eseguiti in un DOM simulato con Testing Library e interrogati come li
userebbe una persona: per ruolo e nome accessibile.

\section{Testing di integrazione}

La fase di testing di integrazione ha lo scopo di verificare la comunicazione
tra i sottosistemi di BiblioFlow. L'attività si concentra sulla catena
completa che va dal route handler al servizio di dominio fino a PostgreSQL,
eseguita contro un database reale e isolato, il cui nome deve dichiarare di
essere di test. Nello specifico verranno esaminati il vincolo
anti-sovrapposizione, le transazioni concorrenti sullo stesso posto, il ciclo
di vita della lista d'attesa e la promozione eseguita dalle automazioni, con
la verifica di notifiche ed eventi prodotti.

\section{Testing di sistema}

Il testing di sistema verrà eseguito secondo un approccio black-box, usando
come confine del sistema le API HTTP esposte dai route handler. La metodologia
di riferimento per la progettazione dei casi di test è il Category Partition
Method, che permette di individuare e definire sistematicamente le
combinazioni di input. I casi progettati sulla versione di partenza ne
fissano il comportamento prima di ogni modifica; la strategia dovrà poi
coprire i cambiamenti apportati al sistema dalle Change Requests.

\section{Testing E2E per la WebApp}

Il testing End-to-End (E2E) ha il compito di validare la WebApp nel suo
complesso, simulando percorsi di utilizzo reali in un browser per garantire la
coerenza tra tutti i sottosistemi. I test si definiscono sui percorsi
principali: l'accesso alle pagine riservate, la prenotazione guidata, il posto
occupato con l'ingresso in lista d'attesa, la promozione dopo una
cancellazione e l'uscita dalla coda. I criteri di successo richiedono che
l'intero ciclo si svolga senza interruzioni sulla build di produzione.
L'intera procedura è automatizzata tramite Playwright.

\chapter{Criteri di Accettazione}

\section{Criteri di Copertura dei Test}

\subsection{Funzionalità incluse nel Testing}

L'attività di testing coprirà le componenti architetturali impattate dalle
modifiche (CR01 e CR02) e le funzionalità core necessarie al loro
funzionamento. Nello specifico:
\begin{itemize}
  \item \textbf{Prenotazione dei posti (CR01):} route handler delle
        prenotazioni, servizio di dominio, autenticazione e middleware, schema
        e vincoli del database.
  \item \textbf{Lista d'attesa e riassegnazione (CR02):} route della coda,
        automazioni periodiche, cancellazioni dello studente e dello staff,
        notifiche ed eventi.
  \item \textbf{Area staff:} route amministrative toccate dalle modifiche e
        controllo dei ruoli su tutte le route dello staff.
  \item \textbf{Interfaccia di prenotazione:} mappa, griglia mobile, messaggio
        di conflitto e gestione della propria coda.
  \item \textbf{Regressione:} accesso, prestiti, statistiche e notifiche
        esistenti, rieseguendo i casi della versione di partenza.
\end{itemize}

\subsection{Funzionalità escluse dal Testing}

Sono escluse dal perimetro di questo piano:
\begin{itemize}
  \item \textbf{Servizi esterni:} il recapito delle email tramite il
        fornitore di posta e l'accesso con Google sono sostituiti da mock;
        resta verificato il comportamento di BiblioFlow quando il servizio
        risponde o fallisce.
  \item \textbf{Monitoraggio e analisi d'uso:} Sentry, Vercel Analytics e
        Speed Insights.
  \item \textbf{Installazione PWA:} il comportamento del service worker e
        della modalità offline, che dipende dal browser e non è toccato dalle
        Change Requests.
\end{itemize}

\section{Criteri e Valori di Accettazione}

Il successo della versione aggiornata di BiblioFlow deve garantire i seguenti
criteri:
\begin{itemize}
  \item \textbf{Copertura del Codice:} almeno il 90\% di statement coverage e
        l'80\% di branch coverage su \file{src/lib} e \file{src/app/api}, e il
        100\% sul middleware. Le soglie sono configurate nello strumento di
        test, che fallisce se non vengono raggiunte.
  \item \textbf{Validazione Funzionale:} il sistema deve rispondere
        correttamente alle modifiche definite nelle Change Requests, in
        particolare:
  \begin{itemize}
    \item \textbf{CA-01, accesso:} ogni operazione su prenotazioni e coda
          richiede la sessione; uno studente non vede né modifica risorse
          altrui e non accede all'area staff;
    \item \textbf{CA-02, concorrenza:} con richieste simultanee sullo stesso
          posto e intervallo viene creata al più una prenotazione, senza
          record parziali;
    \item \textbf{CA-03, lista d'attesa:} l'ordine è quello di arrivo e la
          stessa persona non può avere due richieste attive per lo stesso
          posto e intervallo;
    \item \textbf{CA-04, riassegnazione:} quando il posto si libera il primo
          in coda viene promosso una sola volta, anche con esecuzioni
          concorrenti, e la promozione non confermata scade;
    \item \textbf{CA-05, tracciabilità:} ingresso, uscita, promozione e
          scadenza producono eventi e notifiche correlati;
    \item \textbf{CA-06, regressione:} le funzionalità esistenti non toccate
          dalle modifiche mantengono il loro comportamento.
  \end{itemize}
  \item \textbf{Stabilità:} non devono essere presenti bug bloccanti nei
        processi di testing e la suite E2E deve risultare stabile su
        esecuzioni ripetute.
  \item Tutti i casi di test precedenti classificati come invariati devono
        continuare a passare con successo, garantendo che l'introduzione delle
        modifiche non abbia introdotto regressioni nelle funzionalità
        esistenti.
\end{itemize}

\end{document}
